\section{\S 3.3 盐类水解}\label{sect:盐类水解}

    盐溶液的酸碱性\textcolor{red}{有没有规律？}

    \begin{tabular}{|c|c|}
        \hline
        \textbf{盐溶液} & \textbf{酸碱性} \\
        \hline
        氯化铵、硫酸铝、硝酸铜 & \textcolor{red}{酸性} \\
        \hline
        氯化钠、碘化钾、硝酸钡 & \textcolor{red}{中性} \\
        \hline
        醋酸钠、碳酸钾、次氯酸钙 & \textcolor{red}{碱性} \\
        \hline
    \end{tabular}

    溶液显酸碱性的根本原因是什么？\textcolor{blue}{\cemh{[H+]} 与 \cemh{[OH^-]}}

    \textcolor{red}{这些盐根本不含有 \cemh{[H+]}、\cemh{[OH^-]}， 为什么会造成溶液酸碱性？}
    
    \subsection{盐类水解的原理}\label{subsec:盐类水解的原理}
        
        在盐溶液中盐电离出的离子跟水电离出的 \cemh{H+} 或 \cemh{OH^-} 结合生成\textcolor{red}{\dashuline{\textcolor{black}{{弱电解质}}}}的反应，叫做盐类水解。

        \begin{enumerate}

            \item \textcolor{red}{盐类水解的本质是弱电解质电离的逆过程与水的电离过程相结合；}
            
            \item \textcolor{red}{水解是可逆的，是微弱的（大部分未水解）；}
            
            \item \textcolor{red}{水解可以看成中和反应的逆反应，所以水解是吸热的；}
            
            \[\cemh{\text{盐} + \text{水} <=>[\text{水解}][\text{中和}] \text{酸} + \text{碱}}\quad\mathcolor{blue}{\increment H>0}\]

            \item \textcolor{red}{水解促进了水的电离}
            
            \textcolor{red}{水电离的量看溶液的 \cemh{H+}、\cemh{OH^-} 中大的那个。}

        \end{enumerate} 
        
        \begin{minipage}{0.55\textwidth}
            
            \[
            \begin{matrix}
                \cemh{NH4Cl} & \xleq{} & \cemh{NH4^+} & + \cemh{Cl^-} \\
                & & \rotatebox{90}{\cemh{+}} & \\
                \cemh{H2O} & \cemh{<=>} & \cemh{OH^-} & + \cemh{H+} \\
                & & \rotatebox{90}{\cemh{<=>}} & \\
                & & \cemh{NH3 . H2O} &
            \end{matrix}
            \]
            \vspace{-0.5\baselineskip}

            \begin{center}
                \cemh{H2O <=> H+ + OH^-}
                
                \cemh{NH4+ + OH^- <=> NH3 . H2O}
            \end{center}
        \end{minipage}
        \hfil
        \begin{minipage}{0.35\textwidth}
            \textcolor{blue}{溶液中的 \cemh{OH^-} 被 \cemh{NH4+}}
            
            \textcolor{blue}{结合，\(c(\cemh{OH^-})\)减小，}
            
            \textcolor{blue}{使\(c(\cemh{H+})>c(\cemh{OH^-})\)，}
            
            \textcolor{blue}{溶液显酸性；}
            
            \textcolor{blue}{同时造成水的电离平衡}
            
            \textcolor{blue}{向右移动，}
            
            \textcolor{blue}{水的电离被促进，}
            
            \textcolor{blue}{\(c(\cemh{H+})\)增大。}
        \end{minipage}
        \vspace{0.5\baselineskip}

            总的结果：
            \vspace{-\baselineskip}

            \begin{center}
                \cemh{NH4+ + H2O <=> NH3 . H2O + H+}
                
                \textcolor{red}{盐与水反应，生成了酸和碱。}
            \end{center}

        \vspace{-\baselineskip}
        \[
        \begin{matrix}
            \cemh{CH3COONa} & \xleq{} & \cemh{CH3COO^-} & + \cemh{Na+} \\
            & & \rotatebox{90}{\cemh{+}} & \\
            \cemh{H2O} & \cemh{<=>} & \cemh{H+} & + \cemh{OH^-} \\
            & & \rotatebox{90}{\cemh{<=>}} & \\
            & & \cemh{CH3COOH} &
        \end{matrix}
        \]
        \vspace{-0.5\baselineskip}

        \begin{center}
            \cemh{H2O <=> H+ + OH^-}
            
            \cemh{CH3COO- + H+ <=> CH3COOH}
        \end{center}

        \vspace{-0.5\baselineskip}
        总的结果：
        \vspace{-\baselineskip}

        \begin{center}
            \cemh{CH3COO^- + H2O <=> CH3COOH + OH^-}
            
            \cemh{CH3COONa + H2O <=> CH3COOH + NaOH}
            
            \textcolor{red}{盐与水反应，生成了酸和碱。}
        \end{center}

    \subsection{水解离子方程式的书写}\label{subsec:水解离子方程式的书写}

        \textcolor{red}{书写要点：两个过程叠加的总结果}
        
        书写注意事项：（除双水解反应）

        \begin{enumerate}

            \item 水解是可逆的，用“\cemh{<=>}”表示。
            
            \item 不用“\cemh{v}”、“\cemh{^}”表示\textcolor{red}{（一般水解的程度很小）}
            
            \textcolor{blue}{口诀：书写记三点，水写分子式，中间用可逆，后无沉气析}

            \item 多元弱酸根离子水解必须分步书写。
            
            \cemh{CO3^2- + H2O <=> HCO3^- + OH^-}

            \cemh{HCO3^- + H2O <=> H2CO3 + OH^-}

        \end{enumerate}
        
    \subsection{水解平衡}\label{subsec:水解平衡}

        \subsubsection{水解度}\label{subsubsec:水解度}

            \vspace{-\baselineskip}
            \[
                h = \frac{\text{已水解的离子数}}{\text{总离子数}} \times \SI{100}{\percent} 
            ,\]
            \[
                h = \frac{\text{已水解的离子浓度}}{\text{该离子总浓度}} \times \SI{100}{\percent} 
            .\]
            
            \begin{xmp}
                [水解度例]
                {}
                []
                []
                已知常温下，\SI{0.1}{\mole\per\litre}氯化铵溶液的\(\cemh{pH}=5\)，计算 \cemh{NH4+} 的水解度\(h=\)？
                \[
                    = \num{1.0E-4} = \SI{0.01}{\percent}\qquad\text{\textcolor{red}{“万里挑一”}}
                \]
            \end{xmp}

            影响水解度的因素：
            \(\left\{\begin{matrix}
                \text{\textcolor{red}{加热促进水解（越热越水解）}}\\
                \text{\textcolor{red}{稀释促进水解（越稀越水解）}}\\
                \text{外加物质（依情况具体分析）}
            \end{matrix}\right.\)

        \subsubsection{水解常数}\label{subsubsec:水解常数}
            
            \vspace{-\baselineskip}
            \[\cemh{A- + H2O <=> HA + OH^-}\]
            \[
                K_\text{h} = \frac{\cemh{[HA]}\cdot\cemh{[OH^-]}}{\cemh{[A^-]}} = \frac{\cemh{[HA]}\cdot\cemh{[OH^-]}\cdot\cemh{[H+]}}{\cemh{[A^-]}\cdot\cemh{[H+]}} = \frac{K_\text{w}}{K_\text{a}} \left( = \frac{K_\text{W}}{K_\text{i}}  \right)
            .\]

            \textcolor{red}{\(K_\text{h}\)再一次说明，盐类水解是两个过程共同作用的结果。}

            \textcolor{blue}{根据 \textbf{表~\ref{tab:弱酸-课本3.3}} 的数据可知：大多数弱酸的\(K_\text{a}>\num{1E-7}\)，}

            \textcolor{blue}{则大多数盐水解的\(K_\text{h}<\num{1E-7}\)，}

            \textcolor{red}{\(\therefore\)}\textcolor{blue}{一般来说，盐水解比对应的弱电解质还要弱。}

            如何计算盐溶液的 pH

            \begin{xmp}
                [水解常数计算盐溶液 pH 例]
                {}
                []
                []
                计算\(c=\SI{0.1}{\mole\per\litre}\)醋酸钠溶液的 pH 和电离度\(h\)
                \begin{center}
                    \begin{tabular}{r@{\quad}c@{}c@{}c@{}c@{}c@{}c@{}c}
                        & \(\mathbf{\cemh{CH3COO-}} \) & \(\mathbf{\cemh{+}}\)  & \(\mathbf{\cemh{H2O}}\)  & \(\mathbf{\cemh{<=>}} \) & \(\mathbf{\cemh{CH3COOH}} \) & \(\mathbf{\cemh{+}} \) & \(\mathbf{\cemh{OH^-}} \) \\
                        起始浓度（\si{\mole\per\litre}） & \(\mathbf{0.1}\) &&&& \(\mathbf{0}\) && \(\mathbf{0}\) \\ 
                        平衡浓度（\si{\mole\per\litre}） & \(\mathcolor{blue}{0.1-x}\) &&&& \(\mathcolor{blue}{x}\) && \(\mathcolor{blue}{x}\) \\
                    \end{tabular}
                \end{center}

                \( \displaystyle K_\text{h} = \frac{x^2}{0.1-x}\approx\frac{x^2}{0.1}(x\ll 0.1),\)
                \hfil
                \( \displaystyle K_\text{h} = \frac{K_\text{W}}{K_\text{a}}=\frac{\num{1E-14}}{1.8E-5}=\num{5.56E-10}.\)

                \( \displaystyle \therefore\cemh{[OH^-]}=\sqrt{0.1K_\text{h}}=\sqrt{0.1\times 5.56\times 10^{-10}}=\SI{7.45E-6}{\mole\per\litre} \)

                \( \therefore\cemh{pOH} = 5.13 \quad \therefore\cemh{pH} = 8.87 \) 

                \( \displaystyle h = \frac{\cemh{[OH^-]}}{c} = \frac{\num{7.45E-6}}{0.1} = \num{7.45E-5} \approx \num{1E-4} \)
                \quad \textcolor{red}{“万里挑一”}
            \end{xmp}
            
            \textcolor{blue}{由以上计算过程可以得出：}
            \( \displaystyle \cemh{[OH^-]} = \sqrt{c\cdot K_\text{h}}\)，
            \( \displaystyle h = \sqrt{\frac{K_\text{h}}{c}}\)

            \textcolor{red}{\(\therefore\)盐浓度越大，水解度越小}

            \textcolor{blue}{计算 \cemh{CO3^2-} 和 \cemh{HCO3^-} 的水解常数。}

            \vspace{-\baselineskip}
            \[
                \begin{NiceArray}{r@{\hspace{5em}}l}
                    \cemh{CO3^2- + H2O <=> HCO3^- + OH^-} & \cemh{H2CO3 <=> H+ + HCO3^-} \\[2ex]
                    \cemh{HCO3^- + H2O <=> H2CO3 + OH^-} & \cemh{HCO3^- <=> H+ + CO3^2-}
                \CodeAfter
                    \tikz \draw [red, thick, <->, shorten > = 3pt, shorten < = 2pt] ([yshift=1ex]1-1.base east) -- ([yshift=1ex]2-2.base west);
                    \tikz \draw [red, thick, <->, shorten > = 3pt, shorten < = 2pt] ([yshift=1ex]2-1.base east) -- ([yshift=1ex]1-2.base west);
                \end{NiceArray}
            \]

            \vspace{-0.5\baselineskip}
            \[
                \mathcolor{red}{\therefore}\quad
                K_\text{h1}=\frac{K_\text{W}}{K_\text{a2}}\qquad
                K_\text{h2}=\frac{K_\text{W}}{K_\text{a1}}
            \]

            \textcolor{red}{\(\therefore\)多元弱酸电离，第 1 步\(\gg\)第 2 步}

            \textcolor{red}{\phantom{\(\therefore\)}多元弱酸根离子水解，也是第 1 步\(\gg\)第 2 步}

            \textcolor{blue}{\(\therefore\)碱性：正盐\(>\)酸式盐}

            \subsection{盐类水解的规律}\label{subsec:盐类水解的规律}

                有弱就水解，

                无弱不水解，
                
                越弱越水解，\textcolor{red}{（电离弱，则水解强）（\(K_\text{i}\)越小，则\(K_\text{h}\)越大）}
                
                谁强显谁性。
                
                \begin{table}[htbp]
                    \centering
                    \caption{数据}
                    \label{tab:盐类水解的规律-数据}
                    \begin{tabular}{|c|c|}
                        \toprule
                        \small 弱酸 &\small \( K_\text{a} \) \\
                        \midrule
                        \cemh{HClO} & \( K_\text{a} = \SI{2.95E-8}{} \) \\
                        \midrule
                        \multirow{2}[0]{*}{\cemh{H2S}} & \( K_\text{a1} = \SI{5.7E-8}{} \) \\
                                                        & \( K_\text{a2} = \SI{1.2E-15}{} \) \\
                        \midrule
                        \multirow{3}[0]{*}{\cemh{H3PO4}} & \( K_\text{a1} = \SI{7.52E-3}{} \)  \\
                                                            & \( K_\text{a2} = \SI{6.23E-8}{} \) \\
                                                            & \( K_\text{a3} = \SI{2.2E-13}{} \) \\
                        \bottomrule
                    \end{tabular}
                \end{table}

                \textcolor{blue}{等物质的量浓度的亚硫酸钠溶液和碳酸钠溶液，那个碱性强？}

                \textcolor{blue}{等物质的量浓度的 \cemh{NaHCO3} 溶液和 \cemh{NaClO} 溶液，那个碱性强？}

                强酸强碱盐：不水解
                
                强酸弱碱盐：正离子水解显酸性
                
                弱酸强碱盐：负离子水解显碱性
                
                弱酸弱碱盐：正、负离子水解相互促进，水解程度较大，酸碱性由弱酸弱碱相对强弱决定

                \textcolor{blue}{根据 \textbf{表~\ref{tab:弱酸-课本3.3}} 的数据分析判断，\cemh{NH4F}、\cemh{CH3COONH4}、\cemh{NH4HCO3}、 \cemh{NH4CN} 四种盐溶液的酸碱性。}

            \subsection{酸式盐溶液的酸碱性}\label{subsec:酸式盐溶液的酸碱性}
                
                \textcolor{blue}{\cemh{NaHSO3}、\cemh{NaHCO3}、\cemh{NaHS}、\cemh{NaHC2O4}（草酸氢钠）等酸式盐溶液酸性还是碱性？如何思考判断它们的酸碱性？}

                \subsubsection{多元弱酸的酸式盐}\label{subsubsec:多元弱酸的酸式盐}

                    \vspace{-\baselineskip}

                    \[ \cemh{NaHR \xleq{} Na+ + HR^-} \]

                    \[
                        \left.
                        \begin{array}{ll}
                            \cemh{HR^- <=> H+ + R^2-} & K_{\text{a2}} \\[2ex]
                            \cemh{HR^- + H2O <=> H2R + OH^-} & K_{\text{h2}} = \frac{K_{\text{W}}}{K_{\text{a1}}}
                        \end{array}
                        \right\} \text{\textcolor{red}{对比 \(K_{\text{a2}}\) 与 \(K_{\text{h2}}\) 的大小\footnote{〔数据见\textbf{表~\ref{tab:盐类水解的规律-数据}}〕}}}
                    \]

                    \begin{tabular}{ll}
                        \textcolor{red}{【结论】：}\textcolor{blue}{\cemh{NaHSO3}、\cemh{NaHC2O4} 溶液显酸性；} & \textcolor{red}{（电离大于水解）} \\[1ex]
                        \phantom{\textcolor{red}{【结论】：}}\textcolor{blue}{\cemh{NaHCO3}、\cemh{NaHS} 溶液显碱性。} & \textcolor{red}{（水解大于电离）}
                    \end{tabular}

                \subsubsection{多元强酸的酸式盐\texorpdfstring{\textcolor{red}{（电离显酸性）}}{}}\label{subsubsec:多元强酸的酸式盐}

                    \cemh{NaHSO4 \xleq{} Na+ + H+ + SO4^2-}

                \subsubsection{磷酸盐的酸碱性判断}\label{subsubsec:磷酸盐的酸碱性判断}

                    \begin{xmp}
                        [磷酸盐酸碱性判断例]
                        {}
                        []
                        []
                        根据 \cemh{H3PO4} 的三级电离常数〔见\textbf{表~\ref{tab:盐类水解的规律-数据}}〕判断 \cemh{NaH2PO4} 和 \cemh{Na2HPO4} 两种盐溶液的酸碱性。

                        \begin{enumerate}
                            \item 〔\cemh{NaH2PO4}（\cemh{H2PO4^-}）：〕
                            
                            〔电离：〕\cemh{H2PO4^- <=> H+ + HPO4^2-}

                                \quad \( K_\text{a2} = \num{6.23E-8} \)
                            
                            〔水解：〕\cemh{H2PO4^- + H2O <=> H3PO4 + OH^-}

                                \quad \( K_\text{h3} = \frac{K_\text{W}}{K_\text{a1}} = \frac{\num{1E-14}}{\num{7.52E-3}} = \num{1.33E-12} \)

                            〔\(\because K_\text{a2} > K_\text{h3} \quad \therefore \text{溶液显酸性}\)〕

                            \item 〔\cemh{Na2HPO4}（\cemh{HPO4^2-}）：〕
                            
                            〔电离：〕\cemh{HPO4^2- <=> H+ + PO4^3-}

                                \quad \( K_\text{a3} = \num{2.2E-13} \)
                            
                            〔水解：〕\cemh{HPO4^2- + H2O <=> H2PO4^- + OH^-}

                                \quad \( K_\text{h2} = \frac{K_\text{W}}{K_\text{a2}} = \frac{\num{1E-14}}{\num{6.23E-8}} = \num{1.61E-7} \)

                            〔\(\because K_\text{h2} > K_\text{a3} \quad \therefore \text{溶液显碱性}\)〕

                        \end{enumerate}
                    \end{xmp}

            \subsection{盐类水解的应用}\label{subsec:盐类水解的应用}
            
                \begin{enumerate}
                    \item 判断盐溶液的酸碱性
                    
                    \item 工业上用铝盐或铁盐来净水（水解产生氢氧化铝、氢氧化铁胶体）
                    
                    \cemh{Al^3+ + 3H2O <=> Al(OH)3} (胶体) \cemh{+ 3H+} 
                    
                    \cemh{Fe^3+ + 3H2O <=> Fe(OH)3} (胶体) \cemh{+ 3H+}

                    \textcolor{blue}{\cemh{Al(OH)3} 胶体、\cemh{Fe(OH)3} 胶体具有很强的吸附能力，可吸附水中的杂质并沉降。}

                    \item 配制盐溶液要防止水解
                    
                    \textcolor{blue}{实验室配 \cemh{FeCl3} 溶液是将 \cemh{FeCl3} 溶于稀盐酸中。}

                    \item 用纯碱溶液洗涤油污，为什么加热后效果更好？
                    
                    \cemh{Na2CO3- + H2O <=> NaHCO3 + NaOH} \quad \(\mathcolor{blue}{\increment H>0}\)

                    \item 草木灰不宜与铵态氮肥混合使用
                    
                        \cemh{CO3^2- + H2O <=> HCO3^- + OH^-}

                        \cemh{NH4+ + OH^- <=> NH3 . H2O}
                    
                    \textcolor{blue}{氨水受热分解，使肥效降低。}

                    或：

                        \cemh{CO3^2- + H2O <=> HCO3^- + OH^-}
                        
                        \cemh{NH4+ + H2O <=> NH3 . H2O + H+}

                    \textcolor{blue}{碳酸根离子水解促进了铵离子水解，水解产生的一水合氨增多，氨水受热分解，使肥效降低。}

                    \item 用氯化铵溶液除铁锈
                    
                    \cemh{NH4+ + H2O <=> NH3 . H2O + H+}
                    
                    与 \cemh{Fe2O3} 反应，由于酸性较弱，很难与 \cemh{Fe} 反应。

                    \item 无水氯化铝或无水氯化铁不能从溶液中制取
                    
                    \textcolor{blue}{因为蒸干氯化铝溶液得不到氯化铝}
                    
                    \(\cemh{AlCl3 + 3H2O <=> Al(OH)3 + 3HCl \mathcolor{red}{\uparrow}} \quad \mathcolor{blue}{\increment H>0}\)
                    
                    若再灼烧得到 \cemh{Al2O3}（氢氧化铝受热分解）

                \end{enumerate}

                \noindent\textcolor{red}{【问】加热蒸干硫酸铝溶液得到什么？}

                \textcolor{blue}{硫酸铝晶体}

            \subsection{双水解反应}\label{subsec:双水解反应}

                泡沫灭火器：硫酸铝溶液与碳酸氢钠溶液混合
                    
                \cemh{Al^3+ + 3H2O <=> Al(OH)3 + 3H+}
                    
                \cemh{HCO3^- + H2O <=> H2CO3 + OH^-}
                    
                \textcolor{blue}{两者水解相互促进，平衡向右移动，产生大量 \cemh{Al(OH)3} 和 \cemh{CO2} 气体。}
                    
                总反应：\cemh{Al^3+ + 3HCO3^- \xleq{} Al(OH)3 v + 3CO2 ^}
                
                \noindent 双水解反应——两种离子水解相互促进，生成沉淀或气体，最终完全水解的反应。\footnote{\textcolor{red}{新知识体系下的双水解反应指两种离子水解相互促进的反应，并将原本称作双水解反应的反应称作完全的双水解反应}}
                
                    \textcolor{red}{不是一个水解显酸性、一个水解显碱性就一定会双水解。}
                    
                    \cemh{CH3COO^- + H2O <=> CH3COOH + OH^-}
                    
                    \cemh{NH4+ + H2O <=> NH3 . H2O + H+}

                    \cemh{CH3COO^- + NH4+ + H2O <=> CH3COOH + NH3 . H2O}

                    \textcolor{blue}{\cemh{CH3COO^-} 与 \cemh{NH4+} 水解相互促进，水解程度较大}

                \noindent 哪些离子相互之间会发生双水解反应？

                    正离子：\cemh{Al^3+}、\cemh{Fe^3+}
                    
                    负离子：\cemh{S^2-}（\cemh{HS^-}）、\cemh{CO3^2-}（\cemh{HCO3^-}）、\cemh{SO3^2-}、\cemh{SiO3^2-}、\cemh{ClO^-}
                    
                    \textcolor{red}{注：\cemh{Fe^3+} 与 \cemh{S^2-}（\cemh{HS^-}）发生氧化还原反应、\cemh{Fe^3+} 与 \cemh{SO3^2-} 比较复杂}

                \noindent \textcolor{red}{如何书写双水解反应的方程式}

                    方法一：两个水解的离子方程式叠加
                    
                    方法二：判断最终产物，直接书写配平

                    \cemh{Al^3+ + 3HCO3^-}\quad 〔\cemh{\xleq{} Al(OH)3 v + 3CO2 ^}〕
                    
                    \cemh{Al^3+ + S^2-} \quad 〔\cemh{ + 6H2O \xleq{} 2Al(OH)3 v + 3H2S ^}〕
